% =====================================================================
%  Thème 3 — Colinéarité et parallélisme — Niveau MOYEN — Exercices
% =====================================================================
\documentclass[11pt,a4paper]{article}
\usepackage[utf8]{inputenc}
\usepackage[T1]{fontenc}
\usepackage{lmodern}
\usepackage[french]{babel}
\usepackage{amsmath,amssymb,mathtools}
\usepackage{xcolor}
\usepackage{enumitem}
\usepackage{fancyhdr}
\usepackage[a4paper,margin=2.1cm,headheight=26pt,headsep=10pt]{geometry}

\definecolor{primary}{RGB}{222,70,80}
\definecolor{primarydark}{RGB}{150,30,42}

\pagestyle{fancy}\fancyhf{}
\fancyhead[L]{\small\textcolor{primarydark}{\textbf{Fiche 2} — Calcul vectoriel}}
\fancyhead[R]{\small\textcolor{primarydark}{\textsc{Thème 3 — Moyen — Exercices}}}
\fancyfoot[C]{\small\thepage}
\renewcommand{\headrulewidth}{0.8pt}

\newcommand{\vc}[1]{\overrightarrow{#1}}
\providecommand{\norm}[1]{\left\lVert #1\right\rVert}
\newcommand{\exo}[1]{\medskip\noindent{\color{primary}\bfseries\large Exercice #1.}\par\nopagebreak\smallskip}

\begin{document}
\begin{center}
{\LARGE\bfseries\color{primarydark}Colinéarité et parallélisme — Niveau Moyen}\\[2pt]
{\color{primary}\rule{0.5\linewidth}{1.1pt}}
\end{center}
\vspace{1mm}
{\small\itshape Objectif : remonter d'une condition de colinéarité vers une inconnue, tester
l'alignement avec un paramètre, et décider du parallélisme de deux droites.}
\vspace{2mm}

\exo{1} \textbf{Trouver une composante inconnue.}
Déterminer la (ou les) valeur(s) rendant les deux vecteurs colinéaires.
\begin{enumerate}[label=\textbf{\alph*)},leftmargin=2em,itemsep=2pt]
  \item $\vec u\,(2;3)$ et $\vec v\,(6;y)$.
  \item $\vec u\,(4;-1)$ et $\vec v\,(x;2)$.
  \item $\vec u\,(3;x)$ et $\vec v\,(x;12)$.
\end{enumerate}

\exo{2} \textbf{Alignement avec un paramètre.}
On donne $A\,(1;2)$, $B\,(3;m)$ et $C\,(5;8)$, où $m$ est un réel.
Déterminer $m$ pour que $A$, $B$, $C$ soient alignés.

\exo{3} \textbf{Parallélisme de droites par leurs équations.}
On considère
\[
  d:\ 2x-3y+1=0,\qquad d':\ 4x-6y+7=0,\qquad d'':\ x+y-2=0.
\]
\begin{enumerate}[label=\textbf{\alph*)},leftmargin=2em,itemsep=2pt]
  \item Donner un vecteur directeur de chacune des trois droites.
  \item Quelles droites sont parallèles entre elles ?
  \item Les droites $d$ et $d'$ sont-elles confondues ou strictement parallèles ? Justifier.
\end{enumerate}

\exo{4} \textbf{Parallélisme de droites par des points.}
On donne $A\,(0;1)$, $B\,(2;5)$, $C\,(1;0)$ et $D\,(3;4)$.
\begin{enumerate}[label=\textbf{\alph*)},leftmargin=2em,itemsep=2pt]
  \item Donner un vecteur directeur de $(AB)$ et de $(CD)$.
  \item Les droites $(AB)$ et $(CD)$ sont-elles parallèles ?
  \item Sont-elles confondues (i.e.\ $C$ appartient-il à $(AB)$) ou strictement parallèles ?
\end{enumerate}

\exo{5} \textbf{Construire une droite parallèle.}
Soit $d$ la droite passant par $A\,(2;1)$ et $B\,(5;7)$.
\begin{enumerate}[label=\textbf{\alph*)},leftmargin=2em,itemsep=2pt]
  \item Donner un vecteur directeur de $d$ et sa pente.
  \item Déterminer l'équation réduite de la droite $d'$ parallèle à $d$ et passant par l'origine.
\end{enumerate}

\end{document}
